\section{Introduction}
\label{sec:introduction}

Hybrid pipelines that pair an autoregressive (AR) language model with
a discrete-diffusion language model (DDLM) for chain-of-thought
reasoning have produced contradictory empirical reports.  Some
authors find that an AR planner improves DDLM
denoising~\citep{ye2024diffusionofthought,zhao2025d1};
others find that the same pattern degrades it~\citep{huang2025dcolt};
and reports on whether DDLM sampling diversity is preserved or
collapses under fine-tuning likewise disagree~\citep{liu2025ccdd}.
The picture is muddied further by the fact that recent
DDLMs~\citep{nie2025llada,gong2024diffullama,arriola2025bd3lms} share
a tokenizer family and a semi-AR sampler but are evaluated on
different benchmarks and at different prefix conditions, so the
underlying disagreement is hard to isolate.

We started this project from a narrow observation in our own prior
diagnostic work (E4): on \LLaDA{}-8B-Instruct~\citep{nie2025llada},
prepending a literal \code{"Plan: "} string to a GSM8K
question~\citep{cobbe2021gsm8k} costs $8$\,pp of accuracy versus the
no-prefix baseline ($74\% \rightarrow 66\%$), even though the prefix
contains no content at all.  The damage is format-shape, not plan
quality.  This single number was inconsistent with the simplest
``AR plan helps DDLM'' story and is also inconsistent with the
simplest ``AR plan hurts DDLM'' story: the cost is real, but the
mechanism is not the planner.  It is the prefix surface itself.

This paper argues that the conflicts in the literature collapse once
the failure surface is decomposed into orthogonal axes that can be
measured and intervened on independently.  We identify three.

\begin{enumerate}
    \item \textbf{Interface-format brittleness}: how much the DDLM's
          accuracy drops when the question is wrapped in any
          plan-shaped scaffold, content-free or otherwise.
    \item \textbf{Planner-content trust}: how much the DDLM uses the
          actual content of an upstream AR plan, conditional on the
          prefix shape being absorbed.
    \item \textbf{Sampling-diversity preservation}: whether
          fine-tuning the DDLM collapses, preserves, or expands the
          variety of branches produced by stochastic sampling, on
          which compute-time consensus mechanisms such as
          self-consistency~\citep{wang2023selfconsistency} depend.
\end{enumerate}

On GSM8K-test (N=200) with \LLaDA{}-8B-Instruct and
\Qwen{\{0.5B,1.5B\}}-Instruct~\citep{qwen2024report} planners, we
characterize each axis with a small ($r{=}8$)
LoRA~\citep{hu2022lora,mangrulkar2022peft} and report the following.

First, interface-format brittleness is trainably fixable.  A
prefix-robust LoRA on \LLaDA{} with full FFN module coverage (v3,
22M trainable parameters) lifts no-prefix \ctwo{} from $74\% \to
73.0\%$ (within sampling error) while pulling the empty-prefix
\code{"Plan: "} condition from $66\% \to 74\%$, flattening the
no-prefix-vs-empty-prefix gap from $8$\,pp to within $1$\,pp.  A
partial-coverage variant (v2, 4/7 modules, 10M parameters) achieves
slightly weaker single-shot but stronger branch-vote performance,
exposing a no-prefix-vs-\cmaj{} capacity tradeoff.

Second, planner-content trust is capacity-dependent in opposite
directions across planners.  At v2 capacity, a $1.5$B AR planner
helps ($60\% \to 67\%$, $+7$\,pp) and a $0.5$B planner still hurts
($64\% \to 60\%$).  At v3 capacity this inverts: the $0.5$B
planner improves ($64\% \to 65\%$, $+5$\,pp) and the $1.5$B
planner catastrophically regresses ($60\% \to 54\%$, $-13$\,pp).
Capacity differentially modulates how much the model trusts
upstream content, and the optimal capacity depends on the planner.

Third, format-augmented training expands sampling diversity rather
than collapsing it.  After the prefix-robust LoRA (v2) plus a
late-block commit adapter, the rate at which all $5$ stochastic
branches agree drops from $51.5\%$ to $47.5\%$ on test, and mean
unique answers per problem rises from $1.83$ to $2.07$ ($+13\%$).
\cmaj{} accuracy holds at $81.5\%$ on dev, and \cmajc{} reaches
$82.0\%$ (commit v2) / $82.5\%$ (commit v3) on test.  Diversity
goes up, accuracy is preserved or improved.

Fourth, weight-space distillation of the consensus mechanism
(\cmaj{} into a single forward pass) is design-sensitive in a way
that is easy to miss.  A late-block, answer-span commit adapter
fails to bridge the $6$\,pp gap between greedy \ctwo{} and \cmaj{},
plateauing at \ctwoc{}$=70.5\%$ across two iterations and a
$3.25\times$ capacity bump.  A multi-block, full-response variant
recovers \ctwoc{}$=79.0\%$, within sampling error of the $80\%$
pre-registered target.  Disentangling ablations attribute $+4$\,pp
of the lift to where commit fires in the diffusion schedule and
$+2$\,pp to full-response supervision; full-response training on its
own (with last-block-only commit) adds nothing.

\clearpage
\paragraph{Contributions.}
\begin{enumerate}
    \item \textbf{A three-axis decomposition of hybrid AR/DDLM
          reasoning failure.}  Interface-format brittleness,
          planner-content trust, and sampling-diversity preservation
          are independently measurable on GSM8K with \LLaDA{} and
          \Qwen{\{0.5B,1.5B\}}.  The decomposition reconciles
          conflicting reports in the literature.  Axis 2 is
          characterized as a \emph{phenomenon} over a two-rung
          capacity sweep, not yet a fully-resolved axis on equal
          footing with axes 1 and 3
          (\S\ref{sec:axis1}--\S\ref{sec:axis3},
          \S\ref{sec:discussion}).
    \item \textbf{Trainable interventions on two of the three
          axes.}  Interface-format brittleness is fixable with a
          small ($r{=}8$) prefix-robustness LoRA, with a measurable
          capacity tradeoff: v2 ($4/7$ FFN modules, $10$M
          parameters) and v3 ($7/7$ FFN modules, $22$M parameters)
          trade $+2.5$\,pp on no-prefix \ctwo{} for $-2.0$\,pp on
          \cmaj{} branch-vote, locating a
          single-shot-vs-ensemble Pareto frontier
          (\S\ref{sec:axis1}).  Format-augmented training also
          \emph{expands} sampling diversity rather than collapsing
          it ($5/5$-branch agreement $51.5\% \to 47.5\%$, mean
          unique answers $1.83 \to 2.07$), inverting the canonical
          fine-tuning prediction (\S\ref{sec:axis3}).  A separate
          commit-LoRA distills the consensus mechanism into a
          single forward pass (\S\ref{sec:distill}).
    \item \textbf{Consensus distillation into a DDLM is
          design-sensitive, not architecture-limited, and the
          mechanism is block coverage.}  Late-block answer-span
          surgery plateaus at \ctwoc{}$=70.5\%$ across two
          iterations and a $3.25\times$ capacity bump --- a clean
          negative whose surface reading is ``the adapter cannot
          do this''.  A multi-block ($n{=}3$) full-response
          variant recovers \ctwoc{}$=79.0\%$, within sampling
          error of an $80\%$ pre-registered target.  Disentangling
          ablations attribute most of the lift ($+4$\,pp) to where
          commit fires in the diffusion schedule and $+2$\,pp to
          full-response training; full-response training alone
          (with last-block-only commit) adds $+0.0$\,pp.  Two
          design errors mask each other; block coverage is the
          dominant lever (\S\ref{sec:distill}).
    \item \textbf{\cmaj{} aggregation is structurally lossy by
          $8\text{--}12$\,pp; the failure mode is problem
          comprehension, not arithmetic.}  The oracle ceiling
          exceeds majority-vote accuracy by $8\text{--}12$\,pp
          across four LoRA configurations and three seeds.  The
          gap survives every peer-class verifier we try: $17$
          architectures (text-bag, $0.5$B--$72$B chat-LMs,
          embedding and reward models, process-feature MLPs,
          step-PRMs, symbolic arithmetic) all underperform
          majority vote.  A frontier judge with chain-of-thought
          (Claude Sonnet 4.5~\citep{anthropic2025sonnet}) closes
          $86\%$ of the gap (\cmaj{} $79.1\%$ $\to$ judge
          $85.3\%$, oracle $86.3\%$ on $N{=}500$).  Hand inspection
          of the $15$ \cmaj{}-failed/oracle-recoverable problems
          shows wrong branches contain locally correct arithmetic
          applied to incorrect problem setups, predicting the
          entire verifier leaderboard from a single diagnosis
          (\S\ref{sec:voting_gap}).
\end{enumerate}

Beyond these four headline contributions, additional findings
and honest negatives are surfaced as prose inside the relevant
sections rather than as standalone bullets: a K2 schedule-toggle
ablation locating commit-LoRA as a discrete inference-time toggle
on block-structured mask diffusion (\S\ref{sec:commit_lora_k2}),
an offline-replay per-problem mode router that hits a
pre-registered LOSS threshold (\S\ref{sec:mode_router}), three
parallel aggregation negatives (temporal-SC voting, KL-anchored
schedule-RLHF, LLaDA-$1.5$ base swap;
\S\ref{subsec:parallel_negatives}), a parameter-matched
plain-Qwen-$7$B AR baseline that beats the hybrid \cmajc{} at
fewer active parameters
(\S\ref{subsec:ar_baseline}), a \Qwen{0.5B} self-consistency
rebuttal that rules out generic self-consistency as the
explanation for the diffusion \cmaj{} advantage
(\S\ref{sec:rebuttal}), and the compositionality of param-time
and compute-time consensus (\cmajc{}$=82.5\%$ vs base test
\cmaj{}$=79.0\%$, $+3.5$\,pp; \S\ref{sec:distill}).

\paragraph{Framing: substrate properties, not benchmark wins.}
This paper is about \emph{substrate properties} of current-generation
discrete-diffusion language models for chain-of-thought reasoning
--- what fails, what is fixable, what is structural --- not about
whether a hybrid AR/DDLM stack beats a parameter-matched
monolithic AR baseline on a math benchmark.  At the AR scales we use
as planners, on this benchmark, it does not: a plain
Qwen-$2.5$-$7$B AR baseline reaches $\sim 86.5\%$ on GSM8K-test,
beating our hybrid \cmajc{} of $82.5\%$ at fewer active parameters,
and hybrid accuracy is approximately invariant ($82\text{--}83\%$)
across a $14\times$ planner-scale range from Qwen-$0.5$B to
Qwen-$7$B (\S\ref{subsec:ar_baseline}).  The natural read is that
the substrate-level findings reported here should generalize to
AR/DDLM stacks where the diffusion model is closer to AR
peer-class scale, and that the practical accuracy claim of
\emph{this} hybrid against parameter-matched AR is benchmark- and
scale-bounded.  We take that bet upfront so it does not ambush the
discussion.  Two of the three axes admit trainable interventions
(\S\ref{sec:axis1}, \S\ref{sec:distill}); the third is
characterized as a phenomenon over a two-rung capacity sweep
(\S\ref{sec:axis2}) and we are explicit that a third rung would
upgrade it from a phenomenon to a fully-resolved axis.

\paragraph{Scope and statistical caveats.}
All numbers are measured on GSM8K-test at $N{=}200$, single-seed
$0$ unless otherwise noted, with \LLaDA{}-8B-Instruct at $k{=}64$
denoising steps and a fixed regex answer extractor that tolerates
both \code{Answer:} and \code{\#\#\#\#} formats.  At $N{=}200$,
binomial sampling noise is approximately $\pm 5\text{--}7$\,pp,
which is the right scale to keep in mind when reading the
deltas reported in \S\ref{sec:axis1}--\S\ref{sec:distill}.  The
multi-seed Phase-1 \cmajc{} headline is locked at
$82.2\% \pm 0.29$\,pp ($\sigma$ across $3$ seeds at $N{=}100$),
tighter than the per-seed Wilson interval at $N{=}100$; the other
single-seed deltas in the body should be read as point estimates
and are flagged as such throughout.  The voting-rule gap section
(\S\ref{sec:voting_gap}) additionally uses an $N{=}500$ extension
($200$ frozen indices $+$ $300$ fresh) for the frontier-judge
evaluation; the substrate change moves the base \cmaj{} only from
$79.0\%$ to $79.1\%$, well within sampling noise.
All experiments use a single benchmark (GSM8K-test) and a single
base diffusion model (\LLaDA{}-8B-Instruct); a multi-task
replication (e.g.\ MATH-Easy, ARC-Challenge) and a multi-DDLM
substrate replication (BD3-LMs~\citep{arriola2025bd3lms},
DiffuLLaMA~\citep{gong2024diffullama}) would harden every claim
made here, and we say so plainly in
\S\ref{sec:discussion}-Limitations and \appref{app:open}.

\paragraph{Reproducibility and cost.}
Total compute spend across all four phases of this work is
approximately $\$20.5$ on a single RTX~$4090$ substrate at
\$$0.20$/hr (Runpod), batched into $\sim 100$ pod-hours
(\appref{app:compute}).  All four LoRA adapters (Track 1 v2,
Track 1 v3, Track 2 v2, Track 2 v3) and the three derived datasets
are public on the Hugging Face Hub; the full pipeline
re-runs end-to-end in under a single GPU-day.  Lab-independent
reproducibility on consumer hardware is itself part of the
contribution: the substrate-level findings in this paper do not
require a research-cluster budget to verify, and we view the cost
ceiling as a feature of the design rather than a footnote.
\appref{app:hyperparameters} lists exact training and verifier
configurations; \appref{app:prereg} reproduces the
pre-registration scorecard, including a prediction violated in the
unexpected (good) direction and a post-hoc verifier-aggregation
prediction that hardens \S\ref{sec:voting_gap};
\appref{app:open} flags the remaining open exposures that we did
not close before this draft.
